% =========================================================
% Why Analysis Proofs Feel Different
% =========================================================
\paragraph{Why Analysis Proofs Feel Different.}

Everything in the preceding sections applies to mathematics broadly:
direct proof, induction, contradiction, and the rest are the universal
grammar of mathematical argument. But real analysis --- sequences,
metric spaces, limits, continuity --- introduces a particular
\emph{flavour} of proof work that surprises many students when they
first encounter it.

This section prepares you for that surprise before you arrive at
Volume~III. The goal is not to teach you analysis; it is to give you
the framework so that when you get there, the proofs do not feel alien.

\begin{remark}[The shift from computation to verification]
In calculus, most tasks are \emph{computational}: differentiate this
function, evaluate this integral, find this limit by l'H\^opital's rule.
You apply known procedures and the answer is a number or a function.

In analysis, most tasks are \emph{verificational}: prove that this
sequence converges, show that this space is complete, establish that
this function is continuous. The answer is a logical argument. There is
no procedure to apply; there is a structure to find.

This shift is not just psychological. It changes what ``progress''
means. In a computation, you know you are making progress when the
expressions simplify. In a proof, you know you are making progress
when the definitions are unpacked, the hypotheses are applied, and the
conclusion comes into reach. The habits of mind required are different.
\end{remark}

\begin{remark}[Definitions are not background; they are the machinery]
The single most important habit in analysis is \textbf{returning to
definitions}. Every concept --- convergence, Cauchy, compactness,
continuity --- has a precise definition in terms of quantifiers and
inequalities. That definition is not background knowledge. It is the
working machinery of every proof involving that concept.

Here is the canonical example. To prove that $1/n \to 0$, a calculus
student says: ``obviously it approaches zero.'' An analysis student
unpacks the definition:

\medskip
\noindent\textit{For every $\varepsilon > 0$, there exists $N \in \mathbb{N}$
such that $n \geq N$ implies $|1/n - 0| < \varepsilon$.}

\medskip
\noindent Then they find $N$: choose $N > 1/\varepsilon$ (possible by
the Archimedean property). For $n \geq N$: $1/n \leq 1/N < \varepsilon$.

The proof is three lines, and every line is either the definition
(unpacked), the Archimedean property (a theorem), or arithmetic.
There is no mystery. The entire proof is the definition plus two lines
of calculation. This is the prototype of an analysis proof.

The lesson: when a proof stalls, the first question to ask is always:
\emph{which definition have I not yet unpacked?}
\end{remark}

\begin{remark}[Analysis proofs are structured, not creative]
Here is perhaps the most liberating thing to understand about analysis:
\textbf{proofs are far less creative than they look.}

Experienced analysts do not invent proofs from scratch. They recognize
the type of exercise, recall the standard move for that type, and apply
it. The vocabulary of standard moves is small --- twelve moves cover
virtually all of the proofs in a first course in metric spaces. The
creativity, if any, is in recognizing which move to apply.

This is the same observation that governs the Architecture section of
this chapter (Section~\ref{sec:proof-architecture}), scaled down to the
specific landscape of analysis. The statement map tells you the global
structure. In analysis, a more fine-grained map is available: there are
specific named moves for specific recurring situations. Learning to
name them is the same skill as learning to read a statement and
immediately identify its proof type.

The rest of this section makes that vocabulary explicit.
\end{remark}
